\section{Related Work}
\label{sec:related_work}

\subsection{Quantum Machine Learning and Variational Circuits}
Parameterized quantum circuits (PQCs) represent the predominant approach for implementing machine learning models on near-term NISQ hardware \cite{cerezo2021variational,preskill2018quantum}. 
In variational quantum classifiers, classical data vectors are mapped into quantum Hilbert spaces using non-linear feature maps, such as angle embedding or amplitude embedding \cite{schuld2021machine}. 
The quantum state is subsequently transformed via parameterized single-qubit rotations and multi-qubit entangling operations (e.g., CNOT or CZ gates) before terminal measurement of Pauli observables yields classical expectation values. 
P{\'e}rez-Salinas et al. \cite{perez2020data} introduced the data re-uploading paradigm, demonstrating that interleaving data encoding layers with parameterized unitary operations dramatically expands circuit expressivity without requiring exponential qubit scaling. 
However, training deep variational circuits remains constrained by barren plateaus---regions in the parameter space where cost function gradients vanish exponentially with system size \cite{mcclean2018barren}---and noise-induced barren plateaus arising from environmental decoherence \cite{wang2021noise}.

\subsection{Hybrid Classical-Quantum Architectures in Medical Imaging}
Recent efforts have explored hybrid models combining classical convolutional neural networks (CNNs) with quantum bottlenecks for 2D and 3D medical image segmentation and classification \cite{bergholm2018pennylane}. 
In typical hybrid pipelines, classical encoders compress volumetric inputs down to low-dimensional latent vectors, which are then processed by a quantum circuit before entering a task-specific head. 
A recurring vulnerability across existing literature is the lack of parameter-matched classical controls. 
When a quantum circuit is evaluated without comparison against an equally parameterized classical nonlinear transformation, reported performance gains cannot be rigorously isolated to quantum representation properties rather than arbitrary hyperparameter tuning. 
Furthermore, static fusion mechanisms (such as fixed concatenation or static averaging) fail to account for heteroskedastic noise and varying information density across different patient scans.

\subsection{Multimodal Representation and Feature Modulation}
In neuro-oncology, standard clinical imaging protocols acquire multiple complementary modalities \cite{menze2014multimodal,baid2021rsna}. 
For instance, T1-weighted MRI provides high structural contrast, T2-weighted and FLAIR sequences delineate edema and peritumoural infiltration, and T1-weighted contrast-enhanced (T1ce) scans highlight hypervascular active tumour regions. 
Combining multiple volumetric modalities introduces the challenge of modality interference: as supplementary input channels are incorporated into a fixed-capacity network, redundant or conflicting feature distributions can dilute salient anatomical signals. 
To integrate conditioning signals without unflattening low-dimensional vectors into spatial volumes, Feature-wise Linear Modulation (FiLM) \cite{perez2018film} applies affine transformations directly to intermediate convolutional feature maps:
\begin{equation}
    F' = \gamma(z) \odot F + \beta(z)
\end{equation}
where $\gamma$ and $\beta$ generate channel-wise scaling and shifting vectors based on a conditioning representation $z$. 
FiLM has been successfully applied to visual reasoning and conditional generation; in this work, we demonstrate its utility as a mathematical bridge connecting compact quantum representations with high-resolution 3D convolutional feature maps.
